\chapter{Типизированное вложение причинного ядра в клеточный комплекс}
\label{ch:v10-nonequilibrium-bath-cell-complex-kernel-embedding}

Предыдущий гейт использовал семиузловой путь как контрольный граф. Теперь
проверим, является ли его оператор $A$ самостоятельной матричной догадкой или
следует из цепного комплекса. Возьмём ориентированный путь с пространствами
$C_1\cong\mathbb R^6$, $C_0\cong\mathbb R^7$ и оператором границы
\begin{equation}
 B=\partial_1:C_1\longrightarrow C_0,
 \qquad
 B_{i e}=\begin{cases}-1,&i=e,\\1,&i=e+1,\\0,&\text{иначе}.
 \end{cases}
\end{equation}
Кобграница равна $B^\top$, а нуль-лапласиан и матрица степеней имеют вид
\begin{equation}
 \Delta_0=BB^\top,
 \qquad D=\operatorname{diag}(\Delta_0).
\end{equation}
Тогда типизированный оператор распространения
\begin{equation}
 A_{\rm cell}=\frac12(D-\Delta_0):C_0\longrightarrow C_0
\end{equation}
в точности совпадает с прежней матрицей: его единственные ненулевые элементы
равны $1/2$ на соседних вершинах.

Вложение канонично относительно служебных выборов. Замена ориентаций рёбер
$B\mapsto BS$, где $S^2=I$, не меняет $BB^\top$. Перенумерация вершин
$P$ действует ковариантно:
\begin{equation}
 \Delta_0\mapsto P\Delta_0P^\top,
 \qquad A_{\rm cell}\mapsto PA_{\rm cell}P^\top.
\end{equation}
Матрица $B$ имеет ранг $6$, а $\Delta_0$ --- ранг/ядро $6/1$; единственная
нулевая мода лапласиана представлена постоянной нуль-цепью. Для первых трёх
степеней $A_{\rm cell}$ все элементы вне графового светового конуса снова
тождественно нулевые.

Однако типизация не создаёт диссипативного селектора. Гессиан прежнего
цепного родителя памяти сохраняет ранг/ядро $3/1$ и касательную нулевую моду
$(1,1,3/4,1)$. Кроме того, сам абстрактный комплекс задаёт инцидентность, но
не физическую длину ребра. Размерная карта переменных
$(\tau_{\rm corr},\Delta t,\ell_{\rm edge},v_g)$ по-прежнему имеет
ранг/ядро $2/2$.

\begin{theorem}[Клеточная реализация локального ядра]
Оператор ближайших соседей предыдущего гейта имеет точную типизированную
реализацию через границу и кобграницу ориентированного клеточного пути.
Реализация не зависит от ориентаций рёбер, ковариантна относительно
перенумерации вершин и сохраняет дискретный световой конус.
\end{theorem}

\begin{nogocriterion}[Инцидентность не является метрикой]
Цепной комплекс фиксирует, какие клетки соседствуют, но не выводит физическую
длину ребра, выбор глобального комплекса или коэффициент затухания $r$.
Поэтому типизированное вложение закрывает происхождение локальной поддержки,
но не происхождение абсолютной памяти ванны.
\end{nogocriterion}

\begin{protocol}\begin{enumerate}
\item условная архитектура вложения: \textbf{$10/10$};
\item типизированное происхождение инцидентности: \textbf{$1/1$};
\item происхождение глобального комплекса: \textbf{$0/1$};
\item происхождение физической длины ребра: \textbf{$0/1$};
\item происхождение затухания: \textbf{$0/1$};
\item следующий гейт: \textbf{общий родитель длины ребра клеточного комплекса}.
\end{enumerate}\end{protocol}

Машинный сертификат сохранён в
\path{s2t_v10_cell_birth_four_volume_nonequilibrium_bath_local_causal_propagation_kernel_cell_complex_typed_embedding_gate_results.json}.